\section{Data and Executable Prediction Target}
\label{sec:data}
\subsection{Asset, coverage, and cleaning}
The asset is SPY, denominated in US dollars. The requested data range is 1 January 2010 through 1 January 2026, with the latter endpoint excluded. Yahoo Finance data are obtained using \code{yfinance} with \code{auto_adjust=False}, corporate actions enabled, and automatic repairs disabled \citep{yfinance}. The resulting sample contains \SessionCount{} sessions, from 4 January 2010 through 31 December 2025. Separate raw OHLC, adjusted-close, volume, dividend, and split fields permit independent feature and execution conventions.

Observed dates are checked against the XNYS session calendar, including special closures \citep{exchangecalendars}. There are no missing or extra sessions, duplicate dates, missing fields, or zero-volume observations. The implementation also checks positive finite prices, nonnegative distributions, and valid OHLC ordering. It fails on unsupported data rather than silently filling prices. Extreme observed returns are retained; neither full-sample winsorization nor interpolation is used. Twenty initial observations are excluded for feature warm-up and two terminal observations for unavailable outcomes, leaving 4,002 feature-ready labeled samples across the full history.

SPY has 64 dividend events and no split events in the downloaded interval. The implementation explicitly refuses a nonzero split event pending extended share accounting. Thus the absence of split handling is a stated domain restriction rather than an unobserved modeling error. Raw data remain in a local ignored cache. Provenance records the download options, coverage, package version, retrieval timestamp, and serialized-file digest; published outputs contain derived analysis rather than the full raw price history.

\subsection{Information and target alignment}
Let $t$ index exchange sessions, and let $\mathcal F_t^c$ denote information available after close $t$. A feature vector $x_t$ and forecast $\widehat y_t$ are formed at that time. The position is executed at open $t+1$ and marked through open $t+2$. If $O_t$ is the raw opening price and $D_t$ the distribution with ex-date $t$, the target is
\begin{equation}
 y_t = \frac{O_{t+2}+D_{t+2}}{O_{t+1}}-1.
 \label{eq:target}
\end{equation}
Here the subscript identifies the signal date, not the date on which the outcome becomes observable. A position purchased at the ex-date open is not entitled to that day's distribution; a position carried into the exit ex-date is entitled. Equation~\eqref{eq:target} credits the corresponding receivable at the exit open. Actual payment-date cash timing is abstracted away. We do not add cash dividends to an already total-return-adjusted execution price.

The distinction matters for training. At close $d$, the sample with outcome ending at open $d$ is realized, but a sample ending at open $d+1$ is not. Define its exit session by $e(t)=t+2$. The eligible set is
\begin{equation}
 \mathcal T_d=\operatorname{tail}_{1260}
 \{t:t<d,\ e(t)\le d\}.
 \label{eq:eligible}
\end{equation}
Both the signal-date ordering and label-maturity condition are enforced in code. Training simply on all rows with an earlier feature timestamp would admit an unknown future return.

\subsection{Five trailing features}
Let $A_t$ be adjusted close and $r_t^A=A_t/A_{t-1}-1$. The feature vector contains
\begin{align}
 m_t^{(k)} &= A_t/A_{t-k}-1, \quad k\in\{1,5,20\},\label{eq:mom}\\
 g_t &= \frac{A_t}{\frac1{20}\sum_{j=0}^{19}A_{t-j}}-1,\label{eq:ma}\\
 v_t &= \operatorname{sd}_{19}(r^A_{t-19},\ldots,r^A_t).
 \label{eq:vol}
\end{align}
The notation $\operatorname{sd}_{19}$ indicates the sample standard deviation of 20 returns with denominator 19. All inputs are available in $\mathcal F_t^c$. The return lookbacks represent short and intermediate trends; $g_t$ measures location relative to a recent average, and $v_t$ describes recent uncertainty.

Vendor adjusted histories may incorporate subsequent corporate actions. These features depend on ratios and returns rather than adjusted price levels. A later common multiplicative adjustment on an earlier lookback window cancels from each feature. Appendix~\ref{app:features} states and proves this invariance. Arbitrary vendor revisions need not be common factors and remain a reproducibility limitation, addressed by retaining the input digest.

\section{Regularized Forecasting and Portfolio Construction}
\label{sec:method}
\subsection{Past-only standardization and Ridge fitting}
At each refit date $d$, feature means $\mu_d$ and population standard deviations $s_d$ are computed exclusively on $\mathcal T_d$. Standardized features are $z_{t,d}=(x_t-\mu_d)\oslash s_d$, with the library's constant-feature handling if needed. Ridge solves
\begin{align}
 (\widehat a_d,\widehat\beta_d)
 =\arg\min_{a,\beta}\;&\sum_{t\in\mathcal T_d}
 (y_t-a-z_{t,d}^{\mathsf T}\beta)^2\nonumber\\
 &+\alpha\|\beta\|_2^2.
 \label{eq:ridge}
\end{align}
The intercept is unpenalized. The implementation uses a \code{StandardScaler}/\code{Ridge} pipeline and the deterministic SVD solver \citep{sklearnridge}. Its penalty convention is the sum-of-squares objective in Eq.~\eqref{eq:ridge}, not average squared loss; the numerical scale of $\alpha$ should therefore not be transferred to another normalization without adjustment.

The model is refitted at the first signal and the first decision session of every subsequent calendar month. Between refits, coefficients and scaling remain fixed while $x_t$ changes daily. At least 1,000 eligible observations are required. The actual window contains the most recent 1,260 mature labels, approximately five trading years rather than exactly five calendar years. This gives a controlled amount of estimation history while permitting scheduled adaptation.

Shrinkage is motivated by correlated lookbacks, not by an assumption that any particular predictor must work. A larger penalty limits coefficient magnitudes, but does not eliminate model misspecification or guarantee profitable forecasts. With correlated inputs, individual coefficient signs are conditional descriptions and should not be read as causal estimates of momentum or reversal.

\subsection{Forecasts, positions, and baselines}
Let $d(t)$ be the most recent refit date. The forecast and target position are
\begin{align}
 \widehat y_t &= \widehat a_{d(t)}
 +z_{t,d(t)}^{\mathsf T}\widehat\beta_{d(t)},\\
 w_t &= \ind\{\widehat y_t>0\}.
 \label{eq:position}
\end{align}
Weights are exactly zero or one. There is no leverage, shorting, borrowing, or intermediate fractional allocation. Cash earns zero. The threshold is fixed at zero before results are inspected. This simple map discards forecast magnitude and can trade on small changes near zero; that limitation is part of the tested design, not corrected after seeing the holdout.

Buy-and-hold uses $w_t=1$ throughout the common holding interval. The moving-average baseline uses $w_t^{\mathrm{MA}}=\ind\{g_t>0\}$. Both use the same next-open timing, dividend convention, and fee assumptions as Ridge. The 20-session baseline window is fixed rather than selected from the data.

\subsection{Transaction costs and wealth}
Index the $N$ evaluated holding intervals in chronological order. Set $w_0=0$ before the first entry and let $q_i=|w_i-w_{i-1}|$. For single-side cost $c$, interval growth before terminal liquidation is
\begin{equation}
 G_i(c)=(1-cq_i)(1+w_i y_i).
 \label{eq:fees}
\end{equation}
The final interval is additionally multiplied by $(1-cw_N)$ to liquidate any remaining position. Fees are fractions of pre-trade equity and are deducted before the interval's market return. This is a stated proportional haircut convention, not a detailed broker commission model. With binary weights, no allocation-drift rebalancing is required between position changes.

The main assumption is $c=5/10^4$. Sensitivity scenarios set $c$ to zero and $10/10^4$, holding predictions and positions unchanged. Buy-and-hold pays for its initial purchase and terminal sale. Thus its low turnover is an economic property of the benchmark, not an exemption from cost accounting.

\begin{proposition}[Exact turnover drag]
\label{prop:drag}
For fixed binary weights, fixed $c\in[0,1)$, and positive gross interval wealth, let $K=\sum_{i=1}^N q_i+w_N$. Net and zero-cost terminal wealth satisfy
\begin{equation}
 \frac{V_N(c)}{V_N(0)}=(1-c)^K.
 \label{eq:drag}
\end{equation}
\end{proposition}
Every charged trade contributes the same multiplicative factor and no unchanged position contributes a fee. The full derivation and numerical verification appear in Appendix~\ref{app:accounting}. This identity concerns the frozen path; it is not a model of how trades or liquidity would respond to a different real execution environment.

\subsection{Overfitting controls}
The experiment fixes one asset, five features, one model family, a 1,260-observation window, monthly refits, and a zero forecast threshold. Only $\alpha\in\{0.1,1,10,100\}$ is selected. Validation evaluates the full trading rule after costs rather than optimizing an unrelated in-sample fit metric. The selected setting is locked before final evaluation. This restriction limits search degrees of freedom, although even a four-candidate validation choice can be optimistic.

The pipeline exports each fitted window, label-maturity boundary, scaler, intercept, and coefficient vector. These records provide an operational check of the information set, beyond a general assurance that the backtest is out of sample. Unit tests perturb future prices and confirm that earlier features, fits, and forecasts do not change.
